\SourceLeaf{377}%
threaten the signals and make the observations themselves very difficult. These circumstances prompted the surveyors to form a comparatively large number of fairly small triangles. There are sixteen, and the shortest connecting side is only 18671 metres long. We must observe, however, that heliotropes, which make all signals entirely dispensable and make it just as easy to measure angles accurately in the very largest triangles as in the smallest, were still unknown in Italy at the time.

The angles were measured with eight-inch theodolites made by \textsc{Reichenbach}. The Piedmontese and Austrian officers worked separately, each group determining all the angles of the network independently. This gave two determinations of every angle, and the mean, with due regard to the number of series that had contributed to each, was taken as the final value. The Piedmontese officers' results are mostly based on six series of 10 repetitions each. The Austrian officers' results are mostly based on two series, with some based on three or four. All the measurements are printed in the greatest detail, though without identifying who made each observation.

In a survey of such extent, the accuracy attained in determining the angles is of great interest. The sums of the angles in the individual triangles offer a way to judge it, though in some circumstances this can be somewhat misleading. If the present measurements may be judged by that means, their accuracy is indeed admirable. The largest error in the sum of the angles of any of the 16 triangles is only \(1''16\). The mean error is \(0''70\), and the mean error of individual angles would therefore be only \(0''40\). There is no way to check the results by means of diagonal directions or by forming polygons. Yet comparing the two determinations of each of the 48 angles points quite decidedly to a far greater inaccuracy in the results. In 13 cases we find a difference exceeding \(3''\), and in 5 of those it exceeds \(5''\). Indeed, in one case, in the very first triangle, the Piedmontese officers' determination based on 80 repetitions differs from the Austrian officers' determination based on 48 repetitions by \(9''2\). In the face of such large differences one cannot help suspecting that a proper assessment of the actual accuracy of the measurements also depends on other circumstances of which the work tells us nothing.

